% ============================================================
% METODOLOGÍA — TFG
% Índice Sintético de Turismo Sostenible por CCAA (2016–2024)
% ============================================================

\chapter{Metodología}

\section{Datos y construcción del panel}

Los datos provienen de la plataforma estadística EVASTUR, desarrollada por Dephimática, que centraliza información turística regional procedente de diversas fuentes oficiales. El período analizado abarca de 2016 a 2024. El conjunto inicial contiene 58 indicadores organizados en tres dimensiones: económica (ECO), ambiental (ENV) y social (SOC), referidos a 19 comunidades autónomas españolas más un agregado nacional (ESP).

Todo el pipeline analítico se implementa en R mediante el paquete \texttt{COINr} \citep{becker2022coinr}, que estructura los datos como un objeto \emph{purse} —una colección de \emph{coins} anuales—. Cada \emph{coin} contiene una matriz de datos de dimensión 20 filas (19 CCAA + España) $\times$ número de indicadores por año.

\section{Selección y depuración de indicadores}

\subsection{Criterio 1: Varianza nula}

Se identifican indicadores cuyo rango interanual es nulo en al menos un año, es decir, aquellos para los que $\min = \max$. Este fenómeno causa división por cero en la normalización min-max:

\begin{equation}
    x'_i = 1 + \frac{(x_i - x_{\min})(b - a)}{x_{\max} - x_{\min}}
    \quad \Rightarrow \quad \text{indefinido si } x_{\max} = x_{\min}
\end{equation}

El análisis sobre 9 años identifica 11 indicadores con este problema en al menos un año. Cuando se excluye 2020 (año de disponibilidad limitada), el número se reduce a 5: IAMB0003, IAMB0005, IECO0012, IECO0013 e ISOC0004.

\subsection{Criterio 2: Disponibilidad de datos}

Se calcula el porcentaje de observaciones no ausentes (NA) por indicador:

\begin{equation}
    \text{Disponibilidad}_j = \frac{\text{# observaciones} \, \text{con valor}}{\text{# total de celdas posibles}} \times 100\,\%
\end{equation}

donde el denominador es 19 CCAA $\times$ 9 años = 171. Se descartan indicadores con disponibilidad $< 20\,\%$, criterio que elimina 6 indicadores adicionales: IAMB0004, IAMB0003, IAMB0005, IECO0012, IECO0013 e ISOC0004.

\textbf{Resultado}: panel depurado con $\mathbf{52 \text{ indicadores}}$ activos (58 − 11 + 5 = 52, ya que hay solapamiento entre criterios).

\section{Imputación de valores ausentes mediante MICE}

\subsection{Justificación del método}

El panel contiene valores ausentes de naturaleza heterogénea: algunos indicadores tienen NA sistemáticos en años específicos (ej. cambios en metodología de recopilación), otros presentan patrones aleatorios por disponibilidad de fuentes. La imputación simple por mediana (o eliminación por lista completa) subestima la variabilidad y deja sin cubrir porcentajes altos de NA cuando hay heterogeneidad. Se emplea \textbf{imputación múltiple por ecuaciones encadenadas} (MICE, del paquete \texttt{mice} de R).

\subsection{Especificación técnica}

MICE es un algoritmo iterativo que genera valores imputados mediante:

\begin{enumerate}
    \item Inicialización: imputación provisional por mediana
    \item Iteración $t$ (Gibbs sampling):
    \begin{enumerate}
        \item Para cada variable $X_j$ con NA, ajustar un modelo predictivo sobre todas las demás variables usando las observaciones completas
        \item Predecir los valores ausentes de $X_j$ con el modelo ajustado
        \item Generar valores imputados muestreando de la distribución predictiva condicional
    \end{enumerate}
    \item Repetir el paso 2 hasta convergencia (se fija $\text{maxit} = 5$ iteraciones)
\end{enumerate}

Se emplea el algoritmo \textbf{Predictive Mean Matching} (PMM), que busca entre los 5 vecinos más próximos en el espacio predictor y selecciona uno al azar como valor imputado. PMM garantiza que:
\begin{itemize}
    \item Los valores imputados son reales (preservan el rango empírico)
    \item La distribución marginal de cada variable se preserva
    \item No se introducen artifacts de distribución normal
\end{itemize}

\textbf{Parámetros empleados}:
\begin{itemize}
    \item $m = 1$ (una única cadena imputada)
    \item \texttt{method = "pmm"} (Predictive Mean Matching)
    \item \texttt{maxit = 5} (iteraciones)
    \item \texttt{seed = 42} (reproducibilidad)
\end{itemize}

\textbf{Resultado}: el panel contenía 3.065 NA (32,75\% del total); tras MICE, quedan 118 NA residuales (1,26\%), correspondientes únicamente a IAMB0007, que se imputan por la mediana en fases posteriores.

\section{Tratamiento de valores atípicos}

\subsection{Diagnóstico mediante asimetría y curtosis}

El diagnóstico de normalidad y detección de outliers se realiza mediante dos estadísticos sobre el panel normalizado previo:

\begin{equation}
    \text{Skewness} = \frac{E[(X - \mu)^3]}{\sigma^3}, \quad \text{Kurtosis} = \frac{E[(X - \mu)^4]}{\sigma^4}
\end{equation}

donde $\mu$ es la media y $\sigma$ la desviación típica. Se emplean los criterios del manual de la OCDE \citep{oecd2008handbook}:

\begin{itemize}
    \item \textbf{Asimetría}: $|\text{Skew}| > 2$ indica distribución con cola larga significativa
    \item \textbf{Curtosis}: $\text{Kurt} > 3.5$ indica distribución leptocúrtica (colas pesadas, pico central pronunciado)
\end{itemize}

Un indicador se considera problemático si supera \emph{al menos uno} de estos umbrales.

\subsection{Winsorización}

El tratamiento de valores atípicos se realiza mediante \textbf{winsorización}, implementada por la función \texttt{Treat()} del paquete COINr. La winsorización es un método robusto que no elimina observaciones sino que las reemplaza por el umbral permitido:

\begin{equation}
    x_i^{\text{wins}} = \begin{cases}
        p_L & \text{si } x_i < p_L \\
        x_i & \text{si } p_L \leq x_i \leq p_U \\
        p_U & \text{si } x_i > p_U
    \end{cases}
\end{equation}

donde $p_L$ y $p_U$ son percentiles determinados iterativamente según los umbrales de Skewness y Kurtosis.

\textbf{Parámetros}:
\begin{itemize}
    \item \texttt{f\_t = "winsorise"}: método de tratamiento
    \item \texttt{winmax = 5}: número máximo de observaciones a winsorizar por indicador y año
    \item \texttt{skew\_thresh = 2, kurt\_thresh = 3.5}: umbrales de detección
    \item \texttt{force\_win = TRUE}: fuerza la winsorización incluso si supera \texttt{winmax}
\end{itemize}

El tratamiento se aplica \textbf{año a año}, de modo que cada \texttt{coin} se procesa independientemente. Esto introduce la ventaja de detectar anomalías locales en tiempo, pero pierde información histórica sobre si un valor es extremo solo en ese año o de forma consistente.

\section{Normalización: Min-Max en rango $[1, 100]$}

Se aplica la transformación min-max al rango $[1, 100]$:

\begin{equation}
    x'_i = 1 + (x_i - x_{\min}) \cdot \frac{100 - 1}{x_{\max} - x_{\min}} = 1 + 99 \cdot \frac{x_i - x_{\min}}{x_{\max} - x_{\min}}
    \label{eq:minmax}
\end{equation}

donde $x_{\min}$ y $x_{\max}$ se calculan sobre el panel \emph{completo} (todas las CCAA y todos los años, \texttt{normalise\_by = "all"}), garantizando que:

\begin{itemize}
    \item La escala es comparativa entre años (2016 y 2024 utilizan el mismo rango)
    \item El valor 1 corresponde al peor desempeño observado
    \item El valor 100 corresponde al mejor desempeño observado
    \item Todos los indicadores quedan en la misma escala interpretable
\end{itemize}

\textbf{Justificación del rango $[1, 100]$}: se elige mínimo 1 (no 0) porque la \textbf{media geométrica} empleada en algunos casos de agregación requiere valores estrictamente positivos. El logaritmo natural necesario en la fórmula:

\begin{equation}
    G = \exp\left(\frac{\sum w_j \ln x'_j}{\sum w_j}\right)
\end{equation}

es indefinido en $x = 0$. El desplazamiento de 1 supone esencialmente añadir una constante de regularización que no afecta a la interpretación relativa.

Para indicadores con \texttt{Direction = −1} (donde valores bajos son deseables, ej. emisiones de CO$_2$), la inversión se realiza automáticamente antes de normalizar:

\begin{equation}
    x_{\text{inv}} = -x \quad \Rightarrow \quad \text{min de } x_{\text{inv}} = -\max \text{ del original}
\end{equation}

\section{Análisis de componentes principales (PCA) para ponderación}

\subsection{Estructura jerárquica en dos niveles}

La determinación de ponderaciones empíricas se realiza mediante PCA en estructura jerárquica:

\begin{itemize}
    \item \textbf{Nivel 1}: PCA dentro de cada dimensión (ECO, SOC, ENV) para obtener pesos de indicadores
    \item \textbf{Nivel 2}: PCA entre dimensiones para obtener sus pesos relativos
\end{itemize}

\subsection{Nivel 1: Pesos de indicadores dentro de dimensiones}

Para la dimensión $d \in \{\text{ECO}, \text{SOC}, \text{ENV}\}$, se seleccionan todos los indicadores pertenecientes a $d$ y se aplica PCA sobre la matriz centrada y escalada de esos indicadores:

\begin{equation}
    \mathbf{X}_d \in \mathbb{R}^{N \times k_d}
\end{equation}

donde $N = 19 \text{ CCAA} \times 9 \text{ años} = 171$ observaciones (panel apilado) y $k_d$ es el número de indicadores en la dimensión $d$.

Descomposición de valores singulares:
\begin{equation}
    \mathbf{X}_d = \mathbf{U} \boldsymbol{\Sigma} \mathbf{V}^T
\end{equation}

Los \emph{loadings} del primer componente principal son:

\begin{equation}
    \boldsymbol{\lambda}_1 = \mathbf{v}_1 \cdot \sigma_1
\end{equation}

donde $\mathbf{v}_1$ es el primer vector singular derecho y $\sigma_1$ es el primer valor singular.

\textbf{Pesos normalizados de Nivel 1}:
\begin{equation}
    w_{j}^{(1)} = \frac{|\lambda_{j,1}|}{\sum_{i=1}^{k_d} |\lambda_{i,1}|}
    \label{eq:pesos_l1}
\end{equation}

Se usa el valor absoluto de los loadings porque el signo en PCA es arbitrario (la comparación debe ser sobre magnitud). La normalización asegura $\sum_j w_j^{(1)} = 1$ para cada dimensión.

Se retiene \emph{únicamente} PC1 porque:
\begin{enumerate}
    \item Captura la dirección de máxima varianza compartida
    \item Simplifica la agregación (una única dirección por dimensión)
    \item El uso de PC1 es estándar en construcción de indicadores compuestos \citep{oecd2008handbook}
\end{enumerate}

\textbf{Indicadores con varianza nula} son excluidos previo:
\begin{equation}
    \text{var}(x_j) = \frac{1}{N} \sum_{i=1}^{N} (x_{ij} - \bar{x}_j)^2 \leq 10^{-10} \quad \Rightarrow \quad \text{excluir}
\end{equation}

\subsection{Nivel 2: Pesos de dimensiones}

Con los pesos de Nivel 1, se construye un score dimensional para cada observación:

\begin{equation}
    S_{d,i} = \sum_{j \in d} w_j^{(1)} \cdot x'_{ij}
    \label{eq:score_dim}
\end{equation}

donde $x'_{ij}$ es el valor normalizado del indicador $j$ para la observación $i$.

Sobre los $3 \times N$ scores dimensionales (ENV, ECO, SOC; $N = 171$ observaciones), se aplica PCA nuevamente. Los pesos de Nivel 2 se derivan análogamente:

\begin{equation}
    w_d^{(2)} = \frac{|\lambda_{d,1}|}{\sum_{d'} |\lambda_{d',1}|}
    \label{eq:pesos_l2}
\end{equation}

La estructura jerárquica asegura que:
\begin{itemize}
    \item Cada indicador tiene un peso que refleja su importancia relativa dentro de su dimensión
    \item Cada dimensión tiene un peso que refleja su importancia en el sistema general
    \item Los pesos emergen empíricamente de la estructura de correlaciones, no de juicio subjetivo
\end{itemize}

\section{Métodos de agregación}

Se comparan cinco combinaciones metodológicas para evaluar robustez:

\subsection{Caso 1: Normalización Min-Max + Pesos PCA + Media Aritmética}

Índice para la unidad $u$:

\begin{equation}
    I_u^{(1)} = \sum_{j=1}^{J} w_j \cdot x'_{uj}
    \label{eq:aritmetica}
\end{equation}

donde $w_j$ incluye tanto pesos de Nivel 1 como su agregación por Nivel 2, y $x'_{uj} \in [1, 100]$.

\textbf{Propiedades}:
\begin{itemize}
    \item \textbf{Compensabilidad total}: una región con dimensión muy débil puede compensarse con dimensión muy fuerte
    \item \textbf{Aditividad}: el índice es suma ponderada lineal
    \item \textbf{Rango}: aproximadamente $[1, 100]$ aunque no alcanza los extremos si las ponderaciones no son concentradas
\end{itemize}

\subsection{Caso 2: Normalización Min-Max + Pesos PCA + Media Geométrica}

Índice geométrico ponderado:

\begin{equation}
    I_u^{(2)} = \exp\left(\frac{\sum_{j=1}^{J} w_j \ln x'_{uj}}{\sum_{j=1}^{J} w_j}\right) = \prod_{j=1}^{J} (x'_{uj})^{w_j / \sum_j w_j}
    \label{eq:geometrica}
\end{equation}

De forma equivalente, tomando logaritmos:

\begin{equation}
    \ln I_u^{(2)} = \sum_{j=1}^{J} \tilde{w}_j \ln x'_{uj}
\end{equation}

donde $\tilde{w}_j = w_j / \sum_k w_k$ son pesos normalizados.

\textbf{Propiedades}:
\begin{itemize}
    \item \textbf{No-compensabilidad parcial}: valores muy bajos en cualquier dimensión arrastra el índice de forma sobrelineal
    \item \textbf{Multiplicatividad}: el índice es producto ponderado
    \item \textbf{Interpretación}: penaliza desequilibrios regionales en una óptica de sostenibilidad fuerte (débil complementariedad entre dimensiones)
    \item \textbf{Rango}: $[1, 100]$ teóricamente, aunque rara vez se alcanzan extremos
\end{itemize}

\subsection{Caso 3: Normalización Z-Score + Pesos PCA + Media Aritmética}

Estandarización z-score:

\begin{equation}
    x'_i = \frac{x_i - \bar{x}}{\sigma_x}
    \label{eq:zscore}
\end{equation}

donde $\bar{x}$ y $\sigma_x$ se calculan sobre el panel completo. Produce valores con $E[x'] = 0$ y $\text{Var}(x') = 1$.

Índice con z-score:

\begin{equation}
    I_u^{(3)} = \sum_{j=1}^{J} w_j \cdot x'_{uj}
    \label{eq:case3}
\end{equation}

\textbf{Propiedades}:
\begin{itemize}
    \item Igual a Caso 1 en términos de ranking (preserva el orden relativo)
    \item \textit{Rango}: indefinido, típicamente $[-3, +3]$ pero sin límites conceptuales
    \item Interpretación: desviaciones de la media del panel, no escala absoluta
\end{itemize}

\subsection{Caso 4: Normalización Min-Max + \emph{Benefit of the Doubt} (BoD)}

El \emph{Benefit of the Doubt} es un enfoque derivado del análisis envolvente de datos (DEA) en el que todos los indicadores se tratan como \emph{outputs} (no hay inputs). Para cada unidad $u$, se buscan los mejores pesos $w_j$ que maximicen su índice:

\begin{equation}
    \begin{aligned}
        I_u^{\text{BoD}} &= \max_{w_1, \ldots, w_J} \sum_{j=1}^{J} w_j \, x'_{uj} \\
        \text{s.a.} \quad & \sum_{j=1}^{J} w_j \, x'_{kj} \leq 1 \quad \forall k = 1, \ldots, N \quad (\text{eficiencia factible}) \\
        & w_j \geq 0 \quad \forall j
    \end{aligned}
    \label{eq:bod}
\end{equation}

La restricción de eficiencia asegura que **ninguna otra unidad** puede alcanzar un índice de 1 o superior usando los pesos elegidos para maximizar $u$. Esto equivale a colocar cada unidad en el envolvente de la frontera DEA.

\textbf{Implementación en el contexto de este trabajo}: BoD se aplica sobre los tres scores dimensionales (ENV, ECO, SOC) previamente ponderados con pesos PCA, no sobre los 52 indicadores directamente. Esto reduce la dimensionalidad del problemaBoD y produce resultados interpretables.

\textbf{Propiedades}:
\begin{itemize}
    \item \textbf{Optimización individual}: cada unidad elige sus pesos, sin subjetividad exógena
    \item \textbf{Eficiencia relativa}: el índice refleja posición en la frontera de Pareto
    \item \textbf{Flexibilidad}: puede premiar fortalezas especializadas (ej. si una CCAA es excepcional en una dimensión, BoD lo captura)
    \item \textbf{Limitación}: puede subestimar la importancia de dimensiones equilibradas (sostenibilidad débil vs. fuerte)
\end{itemize}

\section{Estructura informativa del \emph{purse} COINr final}

El \emph{purse} final, una vez completado el pipeline de imputación, tratamiento de outliers y normalización, contiene:

\begin{itemize}
    \item \textbf{9 coins} (años 2016--2024)
    \item \textbf{20 filas} por coin (19 CCAA + España)
    \item \textbf{52 columnas} (indicadores activos)
    \item \textbf{Datasets dentro de cada coin}:
    \begin{itemize}
        \item Raw: datos originales
        \item Imputed: post-MICE
        \item Treated: post-winsorización
        \item Normalised: post-normalización min-max [1--100]
    \end{itemize}
\end{itemize}

Las operaciones de agregación se aplican sobre el dataset \texttt{Normalised} utilizando los pesos derivados del PCA sobre el panel completo apilado.

\end{document}
